
\vspace{-5pt}
\section{Method}
\vspace{-5pt}

We introduce \Method (\method) which enables agents to in-context adapt to new episodes without any re-training.
\method is built using a transformer policy architecture that operates over a long sequence of multi-episode observations and is trained with online RL. 
The novelty of \method is changing the base RL algorithm to more frequently update the policy with increasingly longer contexts within a policy rollout and adding \sinkv to give the model the ability to avoid attending to low-information context.
\Cref{sec:problem} provides the general problem setting of adapting to new episodes.
\Cref{sec:architecture} details the transformer policy architecture.
\Cref{sec:learning} describes the novel update scheme of \method.
Finally, \Cref{sec:impl-details} goes over implementation details.

\subsection{Problem Setting}
\label{sec:problem} 
We study the problem of adaptation to new scenarios in the formalism of meta-RL~\citep{beck2023survey}. 
We have a distribution of training POMDPs $ \mathcal{M}_i \sim p(\mathcal{M})$, where each $ \mathcal{M}_i $ is defined by tuple $ \left( \mathcal{S}_i, \mathcal{S}_i^{0}, \mathcal{O}_i, \mathcal{A}, \mathcal{T}, \gamma, \mathcal{R}_i \right) $ for observations $ \mathcal{O}_i$, states $ \mathcal{S}_i$ which are not revealed to the agent, starting state distribution $ \mathcal{S}_i^{0}$, action space $ \mathcal{A}$, transition function $ \mathcal{T}$, discount factor $ \gamma$, and reward $ \mathcal{R}_i$. In our setting, the states, observations, and reward vary per POMDP, while the action space, discount factor, and transition function is shared between all POMDPs. 

From a starting state $ s_0 \sim \mathcal{S}_i^{0} $, a policy $ \pi$, mapping observations to a distribution over actions, is rolled out for an \emph{episode} which is a sequence of interactions until a maximum number of timesteps, or a stopping criteria. 
We refer to a \emph{trial} as a sequence of episodes within a particular $ \mathcal{M}_i$. The objective is to learn a policy $ \pi $ that maximizes the expected return of an episode.
At test-time the agent is evaluated on a set of holdout POMDPs. 

\subsection{\method Policy Architecture}
\label{sec:architecture} 

Similar to prior work~\citep{grigsby2023amago,team2023human}, \method implements in-context RL via a transformer sequence model that operates over a history of interactions spanning multiple episodes.
At step $ t$ within a trial, \method predicts current action $ a_t$ based on the entire sequence of previous observations $ o_1, \dots , o_t$ which may span multiple episodes.
In the embodied AI settings we study, the observation $ o_t $ consists of an egocentric RGB observation from the robot's head camera along with proprioceptive information and a specification of the current goal.
Each of these observation components are encoded using a separate observation encoding network, and the embeddings are concatenated to form a single observation embedding $ e_t$.
A causal transformer network \citep{vaswani2023attention} $ h_\theta$ inputs the sequence of embeddings $ h_\theta(e_1, \dots, e_t)$. 
From the transformer output, a linear layer then predicts the actions.

The transformer model $ h_\theta$ thus bears the responsibility of in-context learning by leveraging associations between observations within a trial.
This burden especially poses a challenge in our setting of embodied AI since the transformer must attend over a history of thousands of egocentric visual observations. 
Subsequent visual observations are highly correlated, as the agent only takes one action between observations. 
Knowing which observations are relevant to attend to in deciding the current action is thus a challenging problem.
In this work, we build our architecture around full attention transformers using the same architecture as the LLaMA language model~\citep{touvron2023llamaoa}, but modify the number of layers and hidden dimension size to appropriately reduce the parameter count for our setting. 

We also introduce an architectural modification to the transformer called \textbf{\sinkv} to improve the transformer's ability to attend over a long history of visual experience from an embodied agent.
Building off the intuition that learning to attend over a long sequence of visual observations is challenging, we introduce additional flexibility into the core attention operation by prepending the key and value vectors with a per-layer learnable sequence of \emph{sink KV vectors}.
Specifically, for an input sequence $ X \in \mathbb{R}^{n \times d}$ of $ n$ inputs of embedding dimension $ d$, the attention operator projects $ X$ to keys, queries and values notated as $ K, Q, V$ respectively and all elements of $ \mathbb{R}^{n \times d}$ where we assume all hidden dimensions are $d$ for simplicity. 
The standard attention operation computes $ \text{softmax}\left( \frac{Q K^\top}{\sqrt{d}} \right) V $.
We modify calculating the attention scores by introducing learnable vectors $ K_s, V_s \in \mathbb{R}^{s \times d} $ where $ s$ is the specified number of ``sinks''. 
We then prepend $ K_s, V_s$ to the $K, V$ of the input sequence before calculating the attention. 
Note that the output of the attention operation is still $ n \times d$, as in the regular attention operation, as the query vector has no added component.
We repeat this process for each attention layer of the transformer, introducing a new $ K_s, V_s$ in each attention operation.
\sinkv only results in $n_{layers} \times s \times d$ more parameters, which is $0.046\%$ of the 4.5M parameter policy used in this work.

\sinkv gives the sequence model more flexibility on how to attend over the input. 
Prior works observe that due to the softmax in the attention, the model is forced to attend to at least one token from the input \citep{off_by_one_softmax,xiao2023efficient}.
\sinkv removes this requirement by adding learnable vectors to the key and value.
In sequences of embodied visual experiences, this is important as attention heads can avoid attending over any inputs when there is no new visual information in the current observations. 
This flexibility helps the agent operate over longer sequence lengths.

\subsection{\method Learning}
\label{sec:learning} 

\method is updated through online RL, namely PPO~\citep{schulman2017proximal}.
However, for the agent to be able to leverage a long context window for in-context RL, it must also be trained with this long context window.
PPO collects a batch of data for learning by ``rolling out'' the current policy for a sequence of $ T$ interactions in an environment. 
To operate on a long context window spanning an entire trial, the agent must collect a rollout of data that consists of this entire trial.
This is challenging because, in the embodied tasks we consider, we seek to train agents on trials lasting over 64k steps, which consists of at least 130 episodes.  
As typical with PPO, to speed up data collection and increase the update batch size we use multiple environment workers each running a simulation instance that the policy interacts with in parallel. 
With 32 environment workers, this corresponds to $ \approx 130k$ environment steps between every policy update.
PPO policies trained in common embodied AI tasks, such as \obnav, have only $128$ steps between updates and require $ \approx 50k$ updates to converge (for 32 environment workers, 128 steps per worker between updates and 200M environment steps required for convergence)~\citep{yadav2022ovrl}. Executing a similar number of updates would require \method to collect $ \approx 6 $ billion environment interactions. 
\zkn{Need to explain the distinction between update intervals. Would be nice to have a diagram showing entire trial, multiple episodes within it, then which are used for updates (and then show batch axis). \\ Andrew: I think a figure will properly address this, so no need to explain it in the writing.}


\method addresses this problem of sample inefficiency by introducing a \emph{partial update scheme} where the policy is updated multiple times throughout a rollout.
First, at the start of a rollout of length $ T$, all environment workers are reset to the start of a new episode.
Define the number of partial updates as $ K$. 
At step $ i \in  [0, T] $ in the rollout, the policy is operating with a context length of $ i-1$ previous observations to determine the action at step $ i$. 
Every $ T/K$ samples in the rollout, we update the policy.
Therefore, at update $ N$ within the rollout, the agent has collected $ NT/K$ of the $ T$ samples in the rollout. 
The agent is updated using a context window of size $ NT/K$, however, the PPO loss is only applied to the final $ T/K$ outputs.
The policy is changing every $ T/K$ samples in the rollout, so the policy forward pass must be recalculated for the entire $ NT/K$ window rather than caching the previous $ (N-1)T/K$ activations. 
In the last update in the rollout, after collecting the last $T/K$ steps, we update the policy with the loss applied on all steps in the rollout. We refer to this step as \textit{full update}.
At the start of a new rollout, the context window is cleared and the environment workers again reset to new episodes.

\subsection{Implementation Details}
\label{sec:impl-details} 
The transformer is modeled after the LLaMA transformer architecture \citep{touvron2023llamaoa} initialized from scratch. Our policy uses a pretrained visual encoder which is frozen during training. A MLP projects the output of the visual encoder into the transformer.
We only update the parameters of the transformer and projection layers while freezing the visual encoder since prior work shows this is an effective strategy for embodied AI \citep{khandelwal2022simple,vc2023}.
For faster policy data collection, we store the transformer KV cache between rollout steps. 
To fit long context during the training in limited size memory, we used low-precision rollout storage, gradient accumulation \citep{huang2019gpipe} and flash-attention \citep{dao2023flashattention2}.
\asz{What do you mean by ``low-precision rollout"? Clarify what precision you use like ``all models use bfloat16 parameters". And why the distinction with rollout? Is the weight precision different between rollouts and learning? \\Ahmad: I clarified that it's the rollout storage. FYI, I explain the model's precision in the appendix. Mainly, the model is trained with FP32, except for the visual encoder which is FP16. The rollout storage is FP16.}
After each policy update, we shuffle the older episodes in the each sequence and update the KV-cache. 
Shuffling the episode serves as regularization technique since the agent sees the same task for a long time. It also reflects the lack of assumptions about the order of episodes, an episode should provide the same information regardless of whether the agent experiences it at the beginning or at the end of the trial.

We use the VC-1 visual encoder and with the ViT-B size \citep{vc2023}. 
We found the starting VC-1 weights performed poorly at detecting small objects, which is needed for the embodied AI tasks we consider.
We therefore finetuned VC-1 on a small objects classification task.
All baselines use this finetuned version of VC-1.
We provide more details about this VC-1 finetuning in \Cref{sec:visual-encoder-finetuning} and details about all hyperparameters in \Cref{sec:hyperparams}.




